% snippets.tex — Mẫu hình/bảng/code/tham chiếu cho báo cáo.
% Copy đoạn cần dùng, KHÔNG \input cả file này vào báo cáo.

% ---------- Hình ----------
\begin{figure}[H]
  \centering
  \includegraphics[width=0.8\textwidth]{system_topology}
  \caption{Sơ đồ tô-pô hệ thống (nguồn: architecture/system_topology.md).}
  \label{fig:topology}
\end{figure}
% Tham chiếu:  Hình~\ref{fig:topology} mô tả ...

% ---------- Bảng tín hiệu CAN ----------
\begin{table}[H]
  \centering
  \caption{Tín hiệu trong frame 0x100 (nguồn: architecture/can_database.md).}
  \label{tab:can-0x100}
  \begin{tabular}{@{}llll@{}}
    \toprule
    Tín hiệu & Byte & Kiểu & Đơn vị \\
    \midrule
    % TODO[cần đo]: điền theo can_database.md
    speed    & 0--1 & uint16 & km/h \\
    rpm      & 2--3 & uint16 & vòng/phút \\
    \bottomrule
  \end{tabular}
\end{table}

% ---------- Trích code ----------
\begin{lstlisting}[style=c, caption={Hàm encode frame (file:symbol — VD can_encode.c:can_pack_0x100).}, label={lst:can-pack}]
// dán đoạn NGẮN, có ý nghĩa; ghi rõ file:symbol ở caption
\end{lstlisting}

% ---------- Tham chiếu tài liệu ----------
% Trong văn bản:  ... theo chuẩn CAN 2.0B \cite{bosch_can}.
% Trong refs.bib:
% @manual{bosch_can, title={CAN Specification 2.0}, author={Robert Bosch GmbH}, year={1991}}
